\documentclass[11pt, a4paper]{article}

\usepackage[a4paper, margin=2.54cm]{geometry}
\usepackage{fontspec}
\usepackage[spanish]{babel}
\babelfont{rm}{Noto Sans}
\setmonofont{Noto Sans Mono}
\AtBeginDocument{\shorthandoff{"<>}}

\usepackage{xcolor}
\usepackage{listings}
\usepackage[pdftitle={Gramatica Formal del Lenguaje Go},pdfauthor={Flores Pereyra Jairo Jesus}]{hyperref}

\setlength{\emergencystretch}{2.5em}
\def\UrlBreaks{\do\/\do\-\do\.\do\_}

\definecolor{fondocodigo}{HTML}{F8FAFC}
\definecolor{bordecodigo}{HTML}{C7D2FE}
\definecolor{gris}{HTML}{64748B}
\definecolor{kwcolor}{HTML}{2563EB}
\definecolor{ntcolor}{HTML}{0E7490}
\definecolor{opcolor}{HTML}{DB2777}
\definecolor{strcolor}{HTML}{059669}
\definecolor{numcolor}{HTML}{D97706}
\definecolor{puntcolor}{HTML}{475569}
\hypersetup{colorlinks=true,urlcolor=kwcolor}

\lstdefinestyle{gramatica}{
  basicstyle=\ttfamily\renewcommand{\baselinestretch}{1}\fontsize{9}{11.4}\selectfont,
  numbers=none,
  backgroundcolor=\color{fondocodigo},
  frame=leftline,
  rulecolor=\color{bordecodigo},
  framesep=6pt,
  xleftmargin=12pt,
  xrightmargin=6pt,
  aboveskip=8pt,
  belowskip=10pt,
  breaklines=true,
  breakatwhitespace=false,
  breakindent=16pt,
  postbreak=\mbox{\textcolor{gris}{\scriptsize$\hookrightarrow$}\space},
  showstringspaces=false,
  keepspaces=true,
  columns=fullflexible,
  upquote=true,
  alsoletter={_},
  keywordstyle=\color{kwcolor}\bfseries,
  keywordstyle={[2]\color{ntcolor}},
  keywords={package,main,import,type,struct,func,var,int,float64,string,bool,if,else,for,return,true,false},
  morekeywords={[2]inicializacion_opcional,declaraciones_globales,resto_identificador,operador_relacional,declaracion_global,inicializacion_for,lista_expresiones,condicion_simple,lista_sentencias,operador_logico,mas_expresiones,llamada_funcion,declaracion_var,mas_condiciones,literal_struct,mas_parametros,acceso_struct,opcional_else,campos_struct,identificador,tipo_opcional,resto_decimal,parte_decimal,importaciones,actualizacion,resto_entero,resto_cadena,condicional,asignacion,parametros,par_cierra,estructura,expresion,condicion,impresion,sentencia,bucle_for,par_abre,programa,caracter,booleano,funcion,paquete,termino,retorno,simbolo,numero,factor,digito,cadena,bloque,valor,letra,tipo},
  literate={::=}{{\textcolor{opcolor}{\bfseries ::=}}}3
  {:=}{{\textcolor{opcolor}{\bfseries :=}}}2
  {++}{{\textcolor{opcolor}{\bfseries ++}}}2
  {==}{{\textcolor{opcolor}{\bfseries ==}}}2
  {!=}{{\textcolor{opcolor}{\bfseries !=}}}2
  {<=}{{\textcolor{opcolor}{\bfseries <=}}}2
  {>=}{{\textcolor{opcolor}{\bfseries >=}}}2
  {&&}{{\textcolor{opcolor}{\bfseries \&\&}}}2
  {||}{{\textcolor{opcolor}{\bfseries ||}}}2
  {|}{{\textcolor{opcolor}{|}}}1
  {=}{{\textcolor{opcolor}{=}}}1
  {+}{{\textcolor{opcolor}{+}}}1
  {-}{{\textcolor{opcolor}{-}}}1
  {*}{{\textcolor{opcolor}{*}}}1
  {/}{{\textcolor{opcolor}{/}}}1
  {<}{{\textcolor{opcolor}{$<$}}}1
  {>}{{\textcolor{opcolor}{$>$}}}1
  {!}{{\textcolor{opcolor}{!}}}1
  {"}{{\textcolor{strcolor}{"}}}1
  {'}{{\textcolor{strcolor}{'}}}1
  {(}{{\textcolor{puntcolor}{(}}}1
  {)}{{\textcolor{puntcolor}{)}}}1
  {,}{{\textcolor{puntcolor}{,}}}1
  {;}{{\textcolor{puntcolor}{;}}}1
  {.}{{\textcolor{puntcolor}{.}}}1
  {λ}{{\textcolor{numcolor}{λ}}}1
  {->}{{\textcolor{opcolor}{\bfseries$\rightarrow$}}}2
}

\linespread{1.2}
\pagestyle{plain}

\begin{document}

\begin{center}
    {\Large\bfseries Gramática Formal del Lenguaje Go (Golang)}\\[0.6em]
    Univ. Flores Pereyra Jairo Jesus \\
    INF 3631-A: Diseño de Compiladores \\
    Universidad Técnica de Oruro \\[0.4em]
    24 de noviembre de 2026
\end{center}

\noindent\textbf{Enlace al repositorio:}\par
\noindent\raggedright{\small\url{https://github.com/JfloresPereyra/Proyecto-Compilador-Golang-INF-3631A}}

\section{Gramática formal}

\begin{lstlisting}[style=gramatica]
S = <programa>

VN = { <programa>, <paquete>, <importaciones>, <declaraciones_globales>, <declaracion_global>, <estructura>,
<campos_struct>, <funcion>, <parametros>, <mas_parametros>, <tipo_opcional>, <tipo>, <bloque>,
<lista_sentencias>, <sentencia>, <declaracion_var>, <inicializacion_opcional>, <asignacion>, <condicional>,
<opcional_else>, <bucle_for>, <inicializacion_for>, <actualizacion>, <impresion>, <retorno>,
<lista_expresiones>, <mas_expresiones>, <expresion>, <condicion>, <mas_condiciones>, <condicion_simple>,
<operador_logico>, <operador_relacional>, <termino>, <factor>, <valor>, <acceso_struct>, <llamada_funcion>,
<literal_struct>, <par_abre>, <par_cierra>, <identificador>, <resto_identificador>, <numero>, <resto_entero>,
<parte_decimal>, <resto_decimal>, <cadena>, <resto_cadena>, <caracter>, <simbolo>, <booleano>, <letra>,
<digito> }

VT = { package, main, import, "fmt", type, struct, func, var, int, float64, string, bool, if, else, for, return
, fmt.Println, fmt.Printf, true, false, '{', '}', '(', ')', ',', '.', ';', '_', ':', '++', ':=', '=', '+',
'-', '*', '/', '%', '==', '!=', '<', '<=', '>', '>=', '&&', '||', '!', '"', ' ', '\', todas las letras a-z, A
-Z, todos los digitos 0-9 }
\end{lstlisting}

\subsection*{Conjunto de producciones (P)}

\begin{lstlisting}[style=gramatica]
<programa> -> <paquete> <importaciones> <declaraciones_globales>
<paquete> -> package main
<importaciones> -> import "fmt"
<importaciones> -> λ
<declaraciones_globales> -> <declaracion_global> <declaraciones_globales>
<declaraciones_globales> -> λ
<declaracion_global> -> <estructura>
<declaracion_global> -> <funcion>
<estructura> -> type <identificador> struct '{' <campos_struct> '}'
<campos_struct> -> <identificador> <tipo> <campos_struct>
<campos_struct> -> λ
<funcion> -> func <identificador> <par_abre> <parametros> <par_cierra> <tipo_opcional> <bloque>
<parametros> -> <identificador> <tipo> <mas_parametros>
<parametros> -> λ
<mas_parametros> -> ',' <identificador> <tipo> <mas_parametros>
<mas_parametros> -> λ
<tipo_opcional> -> <tipo>
<tipo_opcional> -> λ
<tipo> -> int
<tipo> -> float64
<tipo> -> string
<tipo> -> bool
<tipo> -> <identificador>
<bloque> -> '{' <lista_sentencias> '}'
<lista_sentencias> -> <sentencia> <lista_sentencias>
<lista_sentencias> -> λ
<sentencia> -> <declaracion_var>
<sentencia> -> <asignacion>
<sentencia> -> <condicional>
<sentencia> -> <bucle_for>
<sentencia> -> <impresion>
<sentencia> -> <retorno>
<sentencia> -> <llamada_funcion>
<declaracion_var> -> var <identificador> <tipo> <inicializacion_opcional>
<inicializacion_opcional> -> '=' <expresion>
<inicializacion_opcional> -> λ
<asignacion> -> <identificador> '=' <expresion>
<asignacion> -> <identificador> ':=' <expresion>
<asignacion> -> <acceso_struct> '=' <expresion>
<condicional> -> if <condicion> <bloque> <opcional_else>
<opcional_else> -> else <bloque>
<opcional_else> -> else <condicional>
<opcional_else> -> λ
<bucle_for> -> for <inicializacion_for> ';' <condicion> ';' <actualizacion> <bloque>
<inicializacion_for> -> <asignacion>
<actualizacion> -> <asignacion>
<actualizacion> -> <identificador> '++'
<impresion> -> fmt.Println <par_abre> <lista_expresiones> <par_cierra>
<impresion> -> fmt.Printf <par_abre> <lista_expresiones> <par_cierra>
<lista_expresiones> -> <expresion> <mas_expresiones>
<lista_expresiones> -> λ
<mas_expresiones> -> ',' <expresion> <mas_expresiones>
<mas_expresiones> -> λ
<retorno> -> return <expresion>
<retorno> -> return
<condicion> -> <condicion_simple> <mas_condiciones>
<mas_condiciones> -> <operador_logico> <condicion_simple> <mas_condiciones>
<mas_condiciones> -> λ
<condicion_simple> -> <expresion> <operador_relacional> <expresion>
<condicion_simple> -> '!' <condicion_simple>
<condicion_simple> -> <par_abre> <condicion> <par_cierra>
<condicion_simple> -> <booleano>
<condicion_simple> -> <identificador>
<operador_logico> -> '&&'
<operador_logico> -> '||'
<operador_relacional> -> '=='
<operador_relacional> -> '!='
<operador_relacional> -> '<'
<operador_relacional> -> '<='
<operador_relacional> -> '>'
<operador_relacional> -> '>='
<expresion> -> <termino>
<expresion> -> <expresion> '+' <termino>
<expresion> -> <expresion> '-' <termino>
<termino> -> <factor>
<termino> -> <termino> '*' <factor>
<termino> -> <termino> '/' <factor>
<termino> -> <termino> '%' <factor>
<factor> -> <valor>
<factor> -> <acceso_struct>
<factor> -> <llamada_funcion>
<factor> -> <literal_struct>
<factor> -> <par_abre> <expresion> <par_cierra>
<factor> -> '-' <factor>
<valor> -> <identificador>
<valor> -> <numero>
<valor> -> <cadena>
<valor> -> <booleano>
<acceso_struct> -> <identificador> '.' <identificador>
<llamada_funcion> -> <identificador> <par_abre> <lista_expresiones> <par_cierra>
<literal_struct> -> <identificador> '{' <lista_expresiones> '}'
<par_abre> -> '('
<par_cierra> -> ')'
<booleano> -> true
<booleano> -> false
<cadena> -> " <resto_cadena> "
<resto_cadena> -> <caracter> <resto_cadena>
<resto_cadena> -> λ
<caracter> -> <letra>
<caracter> -> <digito>
<caracter> -> <simbolo>
<simbolo> -> ' '
<simbolo> -> '%'
<simbolo> -> '\'
<simbolo> -> ':'
<simbolo> -> ';'
<simbolo> -> ','
<simbolo> -> '.'
<simbolo> -> '!'
<simbolo> -> '('
<simbolo> -> ')'
<simbolo> -> '+'
<simbolo> -> '-'
<simbolo> -> '*'
<simbolo> -> '/'
<simbolo> -> '='
<simbolo> -> '<'
<simbolo> -> '>'
<simbolo> -> '_'
<identificador> -> <letra> <resto_identificador>
<resto_identificador> -> <letra> <resto_identificador>
<resto_identificador> -> <digito> <resto_identificador>
<resto_identificador> -> '_' <resto_identificador>
<resto_identificador> -> λ
<numero> -> <digito> <resto_entero> <parte_decimal>
<resto_entero> -> <digito> <resto_entero>
<resto_entero> -> λ
<parte_decimal> -> '.' <digito> <resto_decimal>
<parte_decimal> -> λ
<resto_decimal> -> <digito> <resto_decimal>
<resto_decimal> -> λ
<letra> -> a
<letra> -> b
<letra> -> ...
<letra> -> z
<letra> -> A
<letra> -> B
<letra> -> ...
<letra> -> Z
<digito> -> 0
<digito> -> 1
<digito> -> 2
<digito> -> 3
<digito> -> 4
<digito> -> 5
<digito> -> 6
<digito> -> 7
<digito> -> 8
<digito> -> 9
\end{lstlisting}

\end{document}
